
\section{Mechanization, transfer, and trusted boundary}
\label{sec:mechanization}

The previous sections separate the mathematical story into three layers.  First, the fixed
sealed-extension calculus has an exact cubical depth theorem: binary public trace is necessary,
and higher sealing-generated structural trace is computed by open-box filling.  Second, the raw
adequacy bridge turns that theorem into a statement about normalized public signatures and the
minimal trace cost \(\mu\).  Third, the recurrence theorem adds a separate full-coupling and
recent-history hypothesis.

This section explains what the accompanying Cubical Agda development checks and, just as
importantly, what it does not check.  Its purpose is not to replace the paper proofs by a list of
module names.  It is to make the trusted boundary explicit enough that a reader can tell which
claims are internal theorems of the fixed calculus, which claims use actual cubical composition
and partial elements, which claims are surface-language adequacy instances, and which claims
remain semantic transfer principles for future work.

The guiding rule is simple:
\begin{quote}
A checked theorem proves the statement encoded by the checked definitions.  It does not, by
itself, prove that those definitions are an adequate semantics for every external cubical
language or every possible notion of sealed extension.
\end{quote}
That distinction is the reason the paper states its main theorem for the fixed calculus
\(C_{\mathrm{ext}}\), and states broader transfer only through the interface of
Section~\ref{sec:scaling}.

\subsection{The artifact as a theorem-facing development}
\label{subsec:artifact-theorem-facing}

The formalization is organized around the same proof stack as the paper.  There is a raw
extension calculus for payload and trace; an adequacy bridge from surface presentations to raw
candidates; a cubical horn-computational layer; exact depth and recent-history theorems; a
counting layer for sparse and full-coupling recurrences; and a small collection of fixtures and
examples used to audit the intended classification discipline.

It is useful to think of the artifact as having four concentric layers.
\begin{enumerate}
\item The \emph{cubical core layer} checks the open-box and horn-computation arguments using
Cubical Agda's interval, partial elements, composition, filling, transport, paths, and
univalence infrastructure.
\item The \emph{fixed-calculus layer} checks the raw syntax, typing roles, normalized trace
presentations, primitive-versus-derived classification, and \(\mu\)-invariance statements for
\(C_{\mathrm{ext}}\).
\item The \emph{bridge layer} checks that selected surface presentations, especially the modal
fragment SMod, elaborate into the raw calculus while preserving payload, primitive trace,
derived horn trace, historical support, and public trace cost.
\item The \emph{counting and audit layer} checks sparse footprints, the full-coupling envelope,
recurrence arithmetic, refactoring invariance, abstraction barriers, and case-study fixtures.
\end{enumerate}
These layers are related, but they should not be conflated.  A theorem in the fixed-calculus
layer can be postulate-free and still depend on the modeling choice that its record type is the
right representation of a sealed interface.  A theorem in the bridge layer is stronger: it shows
that a specific surface fragment lands in that representation.  A theorem in the cubical core
layer is stronger in a different direction: it uses the actual cubical operations rather than a
synthetic singleton abstraction.

\subsection{Status vocabulary}
\label{subsec:status-vocabulary}

The paper uses the following status vocabulary for artifact claims.  The vocabulary is
intended to prevent a common mistake: reading every checked lemma as if it had the same
semantic scope.

\begin{center}
\small
\setlength{\tabcolsep}{3pt}
\begin{tabular}{p{0.25\linewidth}p{0.66\linewidth}}
\hline
Status & Meaning \\
\hline
\texttt{checked-internal}
& Proved inside the fixed calculus or one of its explicit abstraction records.  This is a
real Agda theorem, but its semantic force is limited to the encoded calculus. \\
\hline
\texttt{checked-actual-cubical}
& Proved using the actual Cubical Agda partial-element, path, composition, filling,
transport, or univalence interface, rather than by assuming a synthetic singleton record. \\
\hline
\texttt{checked-adequacy-instance}
& Proved for a specific surface language or fragment by an explicit elaboration into
\(C_{\mathrm{ext}}\), preserving the trace quantities named in the paper. \\
\hline
\texttt{paper-level-interpretation}
& Explained mathematically in the paper, with a checked counterpart that is slightly more
explicit or more restricted than the prose phrasing.  The finite-basis reading of active-site
cardinality is an example. \\
\hline
\texttt{semantic-assumption}
& A transfer hypothesis about external foundations or languages.  The fixed calculus may
satisfy the hypothesis, but arbitrary external systems are not automatically covered. \\
\hline
\end{tabular}
\end{center}

The absence of local postulates is reported separately from these statuses.  A postulate-free
module is evidence that the encoded theorem was accepted by Agda under the stated options.  It
is not evidence that the encoded object is the only possible, or the most general possible,
semantics of cubical sealed extension.  For example, a postulate-free proof that a synthetic
horn record is contractible would not by itself prove that the actual cubical open-box extension
fiber is contractible.  The development therefore includes a cubical open-box layer whose role
is precisely to connect the theorem-facing structural horn package to Cubical Agda's actual
composition and filling operations.

\subsection{Trusted base and reproducibility}
\label{subsec:trusted-base-reproducibility}

The trusted base is deliberately small but not empty.  It consists of the Agda kernel running
with cubical support and safety options, the pinned Cubical Agda library declared by the
artifact files, the fixed raw extension-calculus specification used as the theorem object, and a
small natural-number/arithmetic layer used by the counting proofs.  The clutching example also
imports the Cubical library's standard higher-inductive circle and suspension interfaces,
together with the usual equivalence, isomorphism, and univalence infrastructure.

The public artifact contains the paper source, theorem-facing Agda modules, the
machine-readable theorem map, audit scripts, and the case-study fixtures.  A reviewer can run
\texttt{scripts/check\_coherence\_depth\_artifact.sh}, or \texttt{make artifact-check} when
GNU \texttt{make} is available.  The script checks the aggregate Agda module, checks the
clutching module separately because of its higher-inductive-type imports, runs the case-study
audit, verifies the paper-to-Agda map, and audits the transitive theorem-facing import closure
for local postulates and local primitive declarations.

The import audit is intentionally mechanical.  It reports the module name, whether local
postulates occur, whether local primitive declarations occur, whether the module is checked
under the intended safety options, and which trusted external imports are used.  This audit is
not a semantic theorem.  It is a reproducibility check for the formal proof boundary.

\begin{remark}[What postulate-free means]
\label{rem:postulate-free-means}
In this section, ``no local postulates'' means that the theorem-facing local modules typecheck
without local \texttt{postulate} declarations or local Agda \texttt{primitive} declarations,
except for the trusted cubical primitives imported from the selected Cubical library.  Auxiliary
files outside the theorem-facing closure may exist, but they are not used by the aggregate
artifact check for the paper's stated theorem map.
\end{remark}

\subsection{Main theorem-to-artifact map}
\label{subsec:theorem-map}

The full theorem map is stored in the artifact as a machine-readable file.  The table below is a
human-readable summary of the load-bearing correspondences.  The purpose is to show where each
mathematical claim enters the checked stack, not to list every helper lemma.

\begin{center}
\scriptsize
\setlength{\tabcolsep}{2pt}
\begin{tabular}{p{0.20\linewidth}p{0.32\linewidth}p{0.30\linewidth}p{0.12\linewidth}}
\hline
Paper claim & Checked components & What is checked & Status \\
\hline
Transparent growth and zero integration latency
& \texttt{Metatheory/InterfaceCalculus.agda}
& Transparent definitions preserve the active interface and contribute no sealed integration
trace. & \texttt{checked-internal} \\
\hline
Raw structural syntax, roles, and normalization
& \texttt{Metatheory/RawStructuralSyntax.agda},
\texttt{RawStructuralTyping.agda}, \texttt{SurfaceNormalizationBridge.agda}
& Raw clauses have the advertised roles; admissible raw extensions elaborate to candidates and
normalize to canonical payload/trace signatures. & \texttt{checked-internal} \\
\hline
Presentation equivalence and \(\mu\)-invariance
& \texttt{CanonicalTelescope.agda}, \texttt{TraceCostNormalForm.agda},
\texttt{PresentationEquivalence.agda}, \texttt{MuInvariance.agda}
& Rebundling, explicit presentation steps, and canonical telescope changes preserve the minimal
public trace cost. & \texttt{checked-internal} \\
\hline
SMod surface adequacy instance
& \texttt{Surface/Modal/*}
& The modal fragment preserves payload, primitive trace, derived horn trace, support, and
\(\mu\) under elaboration into the raw calculus. & \texttt{checked-adequacy-instance} \\
\hline
Structural horn image of raw clauses
& \texttt{SurfaceToHornImage.agda}
& Admissible higher raw structural clauses normalize into the horn-computational image rather
than into primitive higher trace. & \texttt{checked-internal} \\
\hline
Cubical open-box extension theorem
& \texttt{CubicalOpenBox/Base.agda}, \texttt{Extension.agda},
\texttt{Contractible.agda}, \texttt{Substitution.agda}
& The theorem-facing open-box extension package has a canonical point, is contractible, and is
stable under reindexing. & \texttt{checked-actual-cubical} \\
\hline
Horn-to-open-box reduction and Kan subsumption
& \texttt{StructuralHornToOpenBox.agda}, \texttt{KanSubsumption.agda}
& Higher structural obligations reduce to cubical horn/open-box extension fibers; remote-layer
horn obligations are derived. & \texttt{checked-actual-cubical} \\
\hline
Computational replacement
& \texttt{ComputationalReplacement.agda}, \texttt{MuInvariance.agda}
& Explicit higher structural fields are replaced by canonical cubical trace terms and disappear
from \(\mu\)-minimal signatures. & \texttt{checked-internal} with cubical input \\
\hline
Exact stabilization above binary trace
& \texttt{UpperBound.agda}, \texttt{ExactDepth.agda}
& \(\mathrm{Real}_k(X)\) is equivalent to \(\mathrm{Real}_2(X)\) for \(k\geq 2\), with
contractible higher factors in the fixed calculus. & \texttt{checked-internal} \\
\hline
Binary lower bound
& \texttt{AdjunctionBarrier.agda}
& The two-point swap obstruction supplies a concrete witness that unary trace is insufficient
in general. & \texttt{checked-actual-cubical} \\
\hline
Recent-history factorization
& \texttt{ChronologicalWindow.agda}, \texttt{KanSubsumption.agda}
& Surviving primitive support factors through the two most recent normalized public signatures
in the fixed calculus. & \texttt{checked-internal} with cubical input \\
\hline
Sparse and full-coupling recurrences
& \texttt{UniversalRecurrence.agda}, \texttt{SparseDependencyRecurrence.agda},
\texttt{FullCouplingEnvelope.agda}, \texttt{AffineRecurrence.agda}
& Sparse footprints decompose by the active window; the full-coupling endpoint yields the
payload-aware affine recurrence and Fibonacci specialization. & \texttt{checked-internal} \\
\hline
Case-study classification discipline
& \texttt{CaseStudies/*}, \texttt{runs/coherence\_depth\_case\_studies/*},
\texttt{coherence\_depth\_audit.py}
& The intended fixture classifications agree with the payload/trace and sparse/full-coupling
accounting discipline. & checked audit \\
\hline
\end{tabular}
\end{center}

The exact theorem names are deliberately not reproduced in full in the prose.  They are
available in the machine-readable theorem map, and the checker verifies that the paper labels,
Agda module paths, and named Agda declarations exist.  This keeps the paper readable while
still making the mapping auditable.

\subsection{The cubical core: open boxes rather than synthetic singletons}
\label{subsec:cubical-core-artifact}

The most important mechanized boundary is the one between a synthetic horn abstraction and the
actual cubical operations.  Earlier versions of the story could be read as proving that higher
horns are derived because the grammar marked them as derived, or because an abstract singleton
record had been postulated.  The current theorem stack avoids that reading in two ways.

First, the structural syntax still records which fields are intended to be primitive and which
are intended to be horn-computational, but that classification is not the end of the proof.  A
higher structural clause must enter the horn image, reduce to an open-box extension problem,
and receive a canonical cubical replacement.

Second, the open-box package is theorem-facing rather than merely rhetorical.  It contains the
partial boundary, the compatible base/lid data, the filler, and the agreement information needed
for the structural obligation.  The canonical inhabitant is computed from cubical composition
and filling, and the package is stable under substitution.  This is why the paper can say that
higher structural fields are \emph{computed} rather than merely relabeled.

The checked cubical statement should be read in exactly that scope.  It is not a theorem that
all higher filler spaces in Homotopy Type Theory are contractible.  It is a theorem that the
specific theorem-facing structural open-box extension fibers produced by the fixed
sealed-extension calculus have canonical contractible factors.  This is precisely the strength
needed for the public-trace theorem: higher exact factors remain present, but they do not add
primitive public trace.

\begin{theorem}[Checked horn-computational core, artifact reading]
\label{thm:checked-horn-core-artifact-reading}
In the artifact's fixed cubical core, every higher structural clause in the horn image of an
admissible raw declaration is assigned a canonical cubical replacement term, stable under
reindexing, whose public trace presentation has the same normalized boundary data and no
primitive structural field of arity greater than two.
\end{theorem}

\begin{proof}[Artifact proof reading]
The horn-image theorem sends admissible higher raw structural clauses to structural horn
packages.  The horn-to-open-box reduction identifies those packages with theorem-facing cubical
open-box extension fibers.  The open-box theorem supplies the canonical extension and
substitution stability.  The computational-replacement module packages this extension as a
canonical trace presentation, and the \(\mu\)-invariance theorem reads the replacement as
elimination from minimal primitive public trace.
\end{proof}

\subsection{Adequacy instances and the raw bridge}
\label{subsec:adequacy-instances-artifact}

The formalization contains two different kinds of adequacy evidence.

The first is the fixed raw bridge.  This bridge is internal to the paper's theorem object: raw
surface-like clauses are typed, assigned a role, normalized into canonical payload and trace
signatures, and mapped into the horn-computational image when they are higher structural
clauses.  This is the bridge used by the main theorem for \(C_{\mathrm{ext}}\).

The second is a concrete surface-language instance.  The SMod fragment demonstrates that the
adequacy interface is not vacuous.  Its elaboration theorem preserves the quantities the paper
cares about: payload, primitive trace, derived horn trace, support, and \(\mu\).  The explicit
horn presentation, the minimal presentation, and a rebundled presentation all land in the same
normalized public trace cost.

This is a deliberately modest adequacy result.  It does not say that arbitrary Cubical Agda
programs can be parsed and elaborated into \(C_{\mathrm{ext}}\).  It says that one concrete
surface fragment satisfies the bridge conditions, and that the general transfer criterion of
Section~\ref{sec:scaling} identifies exactly what other fragments would have to prove.

\begin{remark}[No parser theorem for arbitrary Cubical Agda]
\label{rem:no-parser-theorem-cubical-agda}
The artifact checks theorem-facing Agda modules and selected surface adequacy instances.  It
does not contain a verified parser, elaborator, or semantic preservation theorem for arbitrary
Cubical Agda code.  Consequently, the paper's claims about external cubical languages should be
read through the transfer interface, not as automatic corollaries of the artifact.
\end{remark}

\subsection{Case-study fixtures}
\label{subsec:case-study-fixtures}

The case-study audit serves a different purpose from the theorem modules.  It checks that the
classification discipline used throughout the paper behaves as advertised on representative
inputs.  In particular, the fixtures separate four phenomena that are easy to confuse:
full-coupling foundational-core growth, sparse local growth, transparent zero-footprint growth,
and user-supplied higher payload.

\begin{center}
\scriptsize
\setlength{\tabcolsep}{2pt}
\begin{tabular}{p{0.24\linewidth}p{0.16\linewidth}p{0.08\linewidth}p{0.24\linewidth}p{0.20\linewidth}}
\hline
Fixture & Regime & \(\mu\) & Trace profile & Boundary demonstrated \\
\hline
\texttt{universe\_extension.yaml}
& full coupling & 2
& unary 1, binary 1, higher horn derived
& fully coupled sealed-extension sequence \\
\hline
\texttt{universe\_extension\_refactored.yaml}
& full coupling & 2
& same normalized trace as the universe fixture
& refactoring preserves public cost \\
\hline
\texttt{global\_modality.yaml}
& full coupling & 3
& unary 1, binary 2, higher horn derived
& global modal payload with full-envelope trace \\
\hline
\texttt{sparse\_datatype.yaml}
& sparse local & 1
& unary 1
& sparse footprint below the full-coupling envelope \\
\hline
\texttt{transparent\_lemma\_extension.yaml}
& transparent zero & 0
& no structural trace obligations
& transparent development lies outside the recurrence \\
\hline
\texttt{promoted\_interface.yaml}
& active-basis totality & 2
& unary 2
& promoted interface package counted by active basis \\
\hline
\texttt{higher\_hit\_payload.yaml}
& payload, not structural trace & 0
& higher constructor payload
& user-supplied higher HIT data is not a structural horn trace \\
\hline
\end{tabular}
\end{center}

The last row is especially important.  The depth-two theorem is not a prohibition on higher
constructors or higher paths.  It is a theorem about sealing-generated structural trace.  A
higher HIT constructor supplied as payload is counted as payload.  A higher horn generated by
structural sealing is counted as derived trace.  The audit fixtures keep that distinction
visible.

\subsection{Abstraction barriers and recurrence proofs}
\label{subsec:barrier-recurrence-artifact}

The recurrence theorem depends on an opacity invariant: once a sealed layer has discharged its
internal obligations, later layers may use only the exported public schema, not the hidden
internal derivation tree that produced it.  This is the formal version of the abstraction
barrier introduced in Section~2.

The artifact checks this discipline in two complementary ways.  The abstraction-barrier modules
verify, for representative later-stage transitions, that inherited obligations are discharged
through the owning layer's opaque record interface rather than by reopening deeper internal
truncation or derivation data.  The decomposition modules then prove the recurrence from an
explicit window decomposition: the next obligation set is split into contributions from the two
most recent sealed layers, and the affine recurrence follows from the cardinalities of those
splits.

This is a useful safeguard.  The recurrence is not proved by first guessing a Fibonacci closed
form and then checking that the numbers fit.  The checked proof goes in the intended direction:
\[
  \text{opaque exported window}
  \quad\Longrightarrow\quad
  \text{finite decomposition of the next footprint}
  \quad\Longrightarrow\quad
  \text{affine recurrence}.
\]
The constant-payload Fibonacci law is then a small arithmetic corollary of the affine
recurrence.

\subsection{What the mechanization establishes}
\label{subsec:what-mechanization-establishes}

Putting the preceding subsections together, the artifact establishes the following theorem
stack for the encoded calculus and checked instances.

\begin{theorem}[Mechanized theorem stack, paper reading]
\label{thm:mechanized-theorem-stack}
Relative to the trusted cubical base and the fixed definitions of \(C_{\mathrm{ext}}\), the
Agda development checks:
\begin{enumerate}
\item the payload/trace and primitive/derived classification discipline for the raw calculus;
\item canonical trace presentations, presentation equivalence, and \(\mu\)-invariance;
\item a concrete SMod adequacy instance preserving payload, trace status, support, and \(\mu\);
\item the structural horn-to-open-box reduction for higher structural clauses;
\item theorem-facing cubical open-box contractibility and substitution stability;
\item computational replacement of explicit higher structural fields;
\item exact stabilization at binary trace and the two-point swap lower bound;
\item recent-history factorization through the two most recent normalized signatures;
\item sparse footprint accounting, the full-coupling affine recurrence, and the constant-payload
Fibonacci specialization;
\item refactoring invariance, abstraction-barrier checks, clutching and lower-bound examples,
and the case-study classification audit.
\end{enumerate}
\end{theorem}

\begin{proof}[Artifact proof reading]
Each item is checked by the theorem-facing modules summarized in
Section~\ref{subsec:theorem-map}.  The aggregate script checks the import closure containing
those modules, the paper-map checker verifies that the paper labels and listed Agda theorem
names exist, and the postulate audit reports no local postulates or local primitive declarations
inside the theorem-facing closure.  The stated result is therefore a theorem of the encoded
artifact relative to the trusted Cubical Agda base and the fixed definitions of the calculus.
\end{proof}

\subsection{What remains outside the mechanized claim}
\label{subsec:outside-mechanized-claim}

The exclusions are as important as the positive claims.

First, the artifact does not prove a broad adequacy theorem for all cubical foundations.  The
paper gives a transfer interface.  A new foundation must still show that its notion of sealed
extension, public trace, horn representation, canonical replacement, and trace cardinality
satisfy that interface.

Second, the artifact does not prove that all higher cubical filler spaces are contractible.  The
contractible factors are the theorem-facing structural open-box extension fibers produced by the
fixed calculus.  Raw higher homotopy, higher inductive payload, and non-structural user data
remain outside the elimination theorem.

Third, the artifact does not prove that ordinary library growth follows the full-coupling
recurrence.  Sparse local declarations and transparent helper definitions are explicitly allowed
and audited.  The recurrence applies only after a two-layer chronological window and full active
basis coupling have been established.

Fourth, the artifact does not remove the trusted base.  The Agda kernel, the Cubical library
interface, and the stated imports remain trusted in the usual way for a mechanized Cubical Agda
formalization.

\subsection{How to read the formalization}
\label{subsec:how-to-read-formalization}

A reader checking the artifact need not begin with the recurrence modules.  The most informative
path follows the paper's proof stack.

Start with the raw extension calculus and canonical trace modules to see how payload, primitive
trace, derived horn trace, and \(\mu\) are represented.  Then inspect the surface-to-horn and
computational-replacement modules; these are the point where ``derived'' becomes a theorem-facing
property rather than a tag.  Next inspect the cubical open-box modules and the horn-to-open-box
reduction; these are the semantic core of the depth-two upper bound.  Then inspect the swap
barrier and exact-depth modules; these supply the lower bound and the final exact-depth theorem.
Finally, inspect the chronological-window and recurrence modules; these explain how the semantic
depth result becomes the two-layer affine law under full coupling.

The case studies should be read last.  They are not the proof of the main theorem.  They are an
audit that the classification vocabulary used by the theorem behaves correctly on small,
representative presentations.

\subsection{Summary}
\label{subsec:mechanization-summary}

The mechanization supports the main message of the paper in the following precise sense.  For
the fixed calculus \(C_{\mathrm{ext}}\), it checks the open-box computation that makes higher
structural trace derived, the replacement theorem that removes such trace from
\(\mu\)-minimal public signatures, the binary lower bound supplied by the swap obstruction, and
the counting consequences under recent-history and full-coupling hypotheses.

The formalization therefore strengthens the paper's central claim, but it does not erase the
semantic boundary.  The fixed theorem is checked.  The SMod bridge is checked as a representative
adequacy instance.  The broad claim that many other cubical foundations or surface languages
instantiate the same interface remains a transfer program: plausible, sharply specified, and
not assumed by the main theorem.
